\documentclass[10pt,a4paper]{amsart}
\usepackage[margin=19mm]{geometry}
\usepackage{amsmath,amssymb,mathtools}
\usepackage[scr=rsfs,cal=euler]{mathalfa}
\usepackage{xfrac}
\usepackage[unicode,hidelinks,bookmarks=false]{hyperref}
\newcommand{\ie}{i.e.\ }
\newcommand{\cf}{cf.\ }
\newcommand{\resp}{respectively\ }
\newcommand{\Subsection}{\subsection}
\providecommand{\nobreakdash}{\nobreak\hbox{-}}
\setlength{\emergencystretch}{2em}
\multlinegap0pt
\usepackage{mathtools}
\usepackage{amssymb}
\usepackage[shortlabels]{enumitem}
\usepackage{tikz}
\usepackage{dsfont}
\usepackage{stackengine}
\newtheorem{prop}{Proposition}
\newtheorem{df}{Definition}
\renewcommand{\P}{\mathds{P}}
\newcommand{\Z}{\mathds{Z}}
\newcommand{\caO}{\mathcal{O}}
\DeclareMathOperator{\rH}{H}
 \DeclareMathOperator{\rd}{rd}
 \DeclareMathOperator{\si}{si}
\title{Height pairing on higher cycles and mixed Hodge structures II}
\author{J.I. Burgos Gil, S. Goswami, G. Pearlstein}
\date{Source: arXiv:2410.17167v3 (CC BY 4.0)}
\begin{document}
\maketitle

\noindent\emph{Source and licence.} Height pairing on higher cycles and mixed Hodge structures II. Version arXiv:2410.17167v3 (CC BY 4.0), distributed by arXiv under the Creative Commons Attribution 4.0 International licence, http://creativecommons.org/licenses/by/4.0/, as recorded in the arXiv licence metadata for this exact version and shown on the abstract page https://arxiv.org/abs/2410.17167v3. Full licence notice: \texttt{licenses/arxiv-2410.17167-CC-BY-4.0.txt}.\par\smallskip

\noindent\textit{Editorial scope.} Counted unit: the article's prop prop:5 (source0.tex lines 3789--3798). The article's complete original proof of it is delivered with it (source0.tex lines 3799--3992, one \texttt{proof} environment of 194 lines). No mathematics is written, repaired or re-derived: the statement and the proof are the article's own lines, in the article's own notation, and no retained line is substituted or re-flowed.  Retained uncounted as the source-local context the proof needs: lines 2867--2882; lines 2927--2932; lines 3453--3465; lines 5368--5380.  Nothing was omitted from any retained span, so the package carries no omission record.  No retained line contains a citation, so the wrapper carries no bibliography and no bibliography entry.  No retained line contains a figure environment, an image-inclusion command or drawing code, so no image is delivered and the package declares no asset dependency.  Editorial formatting only, in addition: an AMS \texttt{amsart} preamble that restates the article's own declarations and macros used by the retained lines, namely the environments \texttt{prop}, \texttt{df} on separate counters, together with the standard packages those lines need.  The retained environments print in this document as Proposition 1, Df 1 in order of appearance, whereas the article prints the same items with the numbering of its own full document.  The complete native source of the article as distributed in the arXiv e-print for this exact version is bundled unchanged as \texttt{sources/167254/source0.tex}.
\begin{prop}\label{prop:6}  Let $W\subset B$ be a smooth
  irreducible variety that has only normal crossing with $D$ (see
  Definition \ref{def:6}). Let
  $\pi\colon \widetilde X\rightarrow X$ be the blow up of $X$ along
  $W$, and $E$ the exceptional divisor of this blow up. By the
  condition on $W$,  $\pi^{-1}(A\cup B)$ is still a normal 
  crossing divisor. Let
  $\widehat A$ and $\widehat B$ be the strict transforms of $A$ and
  $B$, then there is a natural isomorphism of mixed Hodge structures
  \begin{equation}\label{eq:10}
    \varphi \colon \rH^{\ast}(X\smallsetminus A,B)
    \xrightarrow{\cong}
    \rH^{\ast}(\widetilde X\smallsetminus \widehat A,\widehat
    B\cup E).
  \end{equation}
\end{prop}

  \begin{equation}\label{eq:25}
    \pi ^{\ast}\colon
    \Omega ^{\ast}_{X}(\log A\cup B)\otimes \caO(-B)\longrightarrow
    \Omega ^{\ast}_{\widetilde X}(\log \widehat A\cup \widehat B\cup
    E)\otimes \caO(-\widehat B-E).
  \end{equation}

\begin{equation}\label{eq:13}
  \xi_i =
  \begin{cases}
    dz_i/z_i,&\text{ for }i\in I_A\cup I_B,\\
    dz_i,&\text{ for }i\not \in I_A\cup I_B,\\    
  \end{cases}
  \qquad
  \overline{\xi}_i =
  \begin{cases}
    d\bar z_i/\bar z_i,&\text{ for }i\in I_A\cup I_B,\\
    d \bar z_i,&\text{ for }i\not \in I_A\cup I_B.
  \end{cases}
\end{equation}

\begin{df}\label{def:6} Let $X$ be a regular variety of dimension $d$ over a field $K$ of
  characteristic zero.  Let $D$ be a codimension one subvariety  of $X$ and $W$ a
  reduced closed subscheme of $X$. We say that $W$ \emph{has only
    normal crossings with}
  $D$ \emph{at a point} $x$ (or that $D$ has only normal crossings
  with $W$ at $x$) if there is a regular sequence
  $(z_{1},\dots,z_{d})$ in $\caO_{X,x}$   such that the ideal in $\caO_{X,x}$ of each
  component of $D$ containing $x$ is generated by one of the $z_{i}$,
  while the ideal of $W$ in $\caO_{X,x}$ is generated by some of the
  $z_{i}$. We say that $W$ \emph{has only normal crossings with} $D$
  if it has normal crossings with $D$ at every point of $X$. When $W$
  is empty we just say that $D$ has \emph{only normal crossings}.
\end{df}

\begin{prop}\label{prop:5} With the hypothesis and notation of
  Proposition \ref{prop:6},  if $\omega\in E_{X}^\ast(\si A,\rd B)$ then
  \begin{displaymath}
    \pi ^\ast\omega \in E_{\widetilde X}^\ast(\si \widehat A,\rd \widehat B\cup E),
  \end{displaymath}
  Moreover, the isomorphism $\varphi $ of Proposition \ref{prop:6} is
  given for any closed form $\omega \in E_{X}^\ast(\si A,\rd B)$ by
  $\varphi\{\omega\}=\{\pi ^{\ast}\omega\}$, where $\{\omega\}$
  denotes the cohomology class of $\omega$.
\end{prop}

\begin{proof}
  To prove the statement, we work with the local description of the
  blow-up. Let $U\coloneqq X\setminus D$, and
  $V\coloneqq \{z=(z_{1},\cdots , z_{d}); |z_{i}|< \frac{1}{2}\}$ be a
  coordinate neighborhood adapted to $D$ such that there are subsets
  $I_{A}$, $I_{B}$ and 
  $I_{W}$ of $\{1,\dots,d\}$ with 
  \begin{displaymath}
    A\cap V=\{ \prod_{i\in I_{A}} z_{i}=0\},\quad
    B\cap V=\{ \prod_{i\in I_{B}} z_{i}=0\},\quad
    W\cap V=\{ z_{i}=0,\ \forall i\in I_{W}\}.
  \end{displaymath}
  Then $I_{A}$ and $I_{B}$ are disjoint, but $I_{B}\cap I_{W}$ is
  nonempty. The subsets of the shape of $V$ cover $X$ because by
  assumption $W$ has only 
  normal crossings with $D$.

  Let $s=|I_{W}|$. The blow-up
  $\widetilde{V}$ of $V$ at $W\cap V$ is given by the closed subset of
  $V\times \P^{s-1}$ satisfying the set of equations
  $\{z_{i}\alpha_{j}-z_{j}\alpha_{i}=0, \ i,j\in I_{W}\}$,
  where the projective coordinates of $\P^{s-1}$ are
  $(\alpha_{i}\colon)_{i\in I_{W}}$. The blow up $\widetilde V$ can
  be covered by open subsets, $\widetilde V_{j}$, $j\in I_{W}$ defined by the
  condition $\alpha _{j}\not =0$. Fix any element $j\in I_{W}$. We can
  consider the coordinates in $\widetilde V_{j}$ given by
  \begin{displaymath}
    z'_{i}=z_{i}, \text{ for }i\not \in I_{W}, \qquad z'_{i}=\alpha
    _{i}/\alpha _{j}, \text{ for }j\not = i \in I_{W},\qquad z_{j}'=z_{j}.
  \end{displaymath}
  This is not an adapted coordinate neighborhood because the $|z'_{i}|$
  can be bigger than $1/2$. For the sake of brevity we discuss only
  the region $|z_{i}'|\le 1/2$. The remainder of $\widetilde V$ can be
  discussed similarly.
  
  The map $\pi $ is given by the equations
  \begin{equation}\label{eq:12}
    z_{i}=z_{i}', \text{ for }i\not \in I_{W},\qquad 
    z_{i}=z_{i}'z_{j}', \text{ for }i \in I_{W}\smallsetminus \{j\},\qquad
    z_{j}=z_{j}'.
  \end{equation}
  We write $\xi'_{i}$ as in \eqref{eq:13} with respect to the
  variables $z_{i}'$, taking into account that the exceptional divisor
  $E=\{z_{j}=0\}$ is part of the normal crossings divisor.  That is
  \begin{displaymath}
    \xi'_{i}=dz'_{i}/z'_{i},\text{ if }i\in I_{A}\cup
    I_{B}\cup\{j\},\qquad
    \xi'_{i}=dz'_{i},\text{ otherwise.}
  \end{displaymath}
  A simple computation using \eqref{eq:12} shows that 
  \begin{equation}\label{eq:14}
    \pi ^{\ast}\xi_{i}=
    \begin{cases}
      \xi_{i}',&\text{ if }i\not \in I_{W},\\
      \xi'_{j},&\text{ if }i=j\text{ and }j\in I_{A}\cup I_{B},\\
      z_{j}'\xi'_{j} &\text{ if }i=j\text{ and }j\not \in I_{A}\cup I_{B},\\
       \xi_{i}'+\xi_{j}',&\text{ if }i\in (I_{A}\cup I_{B})\cap
                           I_{W}\smallsetminus \{j\},\\
      z_{i}'z_{j}'\xi'_{j}+z_{j}'\xi_{i}'&\text{ if }i\not \in I_{A}\cup
                                           I_{B}\text{ and } i\in  I_{W}\smallsetminus \{j\},
    \end{cases}
  \end{equation}
and similar relations for $\overline{\xi}_{i}$. Since the pullback of
each $\xi_{i}$ is a linear combination of the $\xi_{i}'$  with smooth
coefficients, it is enough to show that, if $f\in E_{X}^0(\si A,\rd
B)$, then
\begin{displaymath}
  \pi ^{\ast}f = f\circ \pi \in E_{\widetilde X}^0(\si \widehat A,\rd \widehat B\cup E).
\end{displaymath}
We start studying the growth condition of the function and later we
will study the growth of the derivatives. 
  Now, let $f\in E^{0}(\si A, \rd B)$. There exists an integer $N_1\in
  \Z$ and for each integer $N_2\in \Z$  there is a
  constant $C_{N_1,N_2}$
\begin{displaymath}
 \left | f
    \right|\le C_{N_1,N_2}
    \prod_{i\in I_A}|\log|z_i||^{N_1}
    \prod_{k\in I_B}|\log|z_k||^{N_2}.
  \end{displaymath}
  Then 
  \begin{align*}
    \left | f\circ \pi 
    \right|&\le C_{N_1,N_2}
    \prod_{i\in I_A\cap (I_{W}\smallsetminus \{j\})}|\log|z'_{j}z'_i||^{N_1}
     \prod_{i\in I_A\smallsetminus (I_{W}\smallsetminus \{j\})}|\log|z'_i||^{N_1}\\
    &\qquad \qquad \qquad \prod_{i\in I_B\cap (I_{W}\smallsetminus \{j\})}|\log|z'_{j}z'_i||^{N_2}
     \prod_{i\in I_B\smallsetminus (I_{W}\smallsetminus \{j\})}|\log|z'_i||^{N_2}
  \end{align*}
  Since $|z'_{i}|\le 1/2$ for all $i$, there is a constant $C>0$ such
  that
  \begin{displaymath}
    \log|z'_{j}z_i'|=\log|z'_{j}|+\log|z'_i|\le C\log|z'_{j}|\log|z'_i|. 
  \end{displaymath}
  Therefore, there is another constant $C'_{N_{1},N_{2}}$ such that 
  \begin{equation}\label{eq:20}
    \left | \pi ^{\ast }f 
    \right|\le C'_{N_1,N_2} |\log|z'_j||^{aN_{1}+bN_2}
    \prod_{i\in I_A\smallsetminus \{j\}}|\log|z'_i||^{N_1}
   \prod_{i\in I_{B}\smallsetminus \{j\}}|\log|z'_i||^{N_2},
 \end{equation}
 where $a=|I_{A}\cap I_{W}|$ and $b=|I_{B}\cap I_{W}|$. By hypothesis
 $W\subset B$ so $I_{B}\cap I_{W}\not = \emptyset $ and $b>0$. Since
 $N_{1}$ is fixed, and $N_{2}$ is arbitrary, we can make
 $aN_{1}+bN_{2}$ arbitrarily negative. This shows that   $\pi
 ^{\ast}f$ satisfies the required bounds.

 The chain rule together with equations \eqref{eq:12} show that
 \begin{displaymath}
   \partial/\partial z'_{i}=
   \begin{cases}
     \partial/\partial z_{i},
     &\text{ if }i\not \in I_{W},\\
     z_{j}\partial/\partial z_{i},
     &\text{ if }i\in I_{W}\smallsetminus \{j\}\\
     \sum_{k\in I_{W}} (z_{k}/z_{j})\partial/\partial z_{k},
     & \text{ if }i = j. 
   \end{cases}
 \end{displaymath}
 This implies that the differential operator  $\partial'{} ^{\alpha,\beta  }$
 can be written as a linear combination of terms of the form
 \begin{math}
   z^{\alpha '',\beta ''}\partial^{\alpha ',\beta '}
 \end{math}
 where
 \begin{equation}
   \label{eq:17}
   \begin{gathered}
      \alpha '_{j}\le \alpha _{j},\qquad \alpha '_{i}\ge \alpha _{i}
 \text{ for } i\in I_{W}\smallsetminus \{j\},\qquad \alpha '_{i}=\alpha _{i}\text{ for }i\not \in
 I_{W},\\
 \alpha _{i}''=\alpha '_{i}-\alpha _{i}, \qquad
 \alpha ''_{j}=\sum _{i\in I_{W}\smallsetminus \{j\}}\alpha _{i}+\alpha '_{j}-\alpha _{j},
\\
      \beta '_{j}\le \beta _{j},\qquad \beta '_{i}\ge \beta _{i}
 \text{ for } i\in I_{W}\smallsetminus \{j\},\qquad \beta '_{i}=\beta _{i}\text{ for }i\not \in
 I_{W},\\
 \beta _{i}''=\beta '_{i}-\beta _{i}, \qquad
 \beta ''_{j}=\sum _{i\in I_{W}\smallsetminus \{j\}}\beta _{i}+\beta '_{j}-\beta _{j},
   \end{gathered}
 \end{equation}
 Thus, we have to show that any function of the form
 $z^{\alpha '',\beta ''}\partial^{\alpha ',\beta '}f$ with
 $\alpha', \alpha '', \beta ', \beta ''$ satisfying the conditions
 \eqref{eq:17}, when expressed in terms of the coordinates $z'$
 satisfies the bounds required to $\partial'{}^{\alpha ,\beta }\pi
 ^{\ast}f$. For brevity of the exposition we write only the case
 $\partial^{\alpha }$ as the general case follows the same principle
 and is only more cumbersome to write. We have
 \begin{displaymath}
   \left |z^{\alpha ''} \partial ^{\alpha '}f\right | \le C_{N_1,N_2,\alpha' }
    \frac{\prod_{i\in I_A}|\log|z_i||^{N_1}
      \prod_{j\in I_B}|\log|z_j||^{N_2}}
    {\prod_{i\in I_{A}\cup I_{B}}|z_{i}|^{\alpha' _{i}-\alpha _{i}''}}.
  \end{displaymath}
  Using the relations \eqref{eq:17} and \eqref{eq:12} one proves
  that,
  \begin{displaymath}
    \prod_{i\in I_{A}\cup I_{B}}|z_{i}|^{\alpha' _{i}-\alpha _{i}''}=
    |z'_{j}|^{\alpha''' _{j}}\prod_{i\in I_{A}\cup I_{B}\smallsetminus \{j\}}|z'_{i}|^{\alpha_{i}},
  \end{displaymath}
where, 
\begin{displaymath}
  \alpha ^{'''}_{j}=
  \begin{cases}
    \alpha _{j}-\sum_{i\in I_{W}\smallsetminus (I_{A}\cup
    I_{B})}\alpha _{i}, &\text{ if }j\in I_{A}\cup I_{B}\\
    \alpha _{j}-\sum_{i\in I_{W}\smallsetminus (I_{A}\cup
    I_{B}\cup \{j\})}\alpha _{i}-\alpha '_{j}, &\text{ if }j\not \in
                                                 I_{A}\cup I_{B}.
  \end{cases}
\end{displaymath}
In any case, $\alpha '''_{j}\le \alpha _{j}$.  Combining the bound
 \eqref{eq:20} that we obtained for the numerator, and the previous
 computation taking into account that $\alpha _{j}'''\le \alpha _{j}$
 we deduce
 \begin{displaymath}
   \left |\partial'{} ^{\alpha }\pi ^{\ast }f\right |\le
   C'_{N_{1},N_{2},\alpha }|\log|z'_j||^{aN_{1}+bN_2}\frac 
    {\displaystyle  \prod_{i\in I_A\smallsetminus \{j\}}|\log|z'_i||^{N_1}
   \prod_{i\in I_{B}\smallsetminus
     \{j\}}|\log|z'_i||^{N_2}}{\displaystyle  \prod_{i\in I_{A}\cup
     I_{B}\cup\{j\}}|z_{i}|^{\alpha _{i}}} 
 \end{displaymath}
 with $b> 0$


The last statement follows from the fact that the map
 \begin{displaymath}
 \pi ^{\ast}\colon E_{X}^\ast(\si A,\rd B) \to  E_{\widetilde X}^\ast(\si \widehat A,\rd
 \widehat B\cup E)
 \end{displaymath}
  is compatible with the map \eqref{eq:25}.
\end{proof}

\end{document}
